% !TEX root = ../main.tex

\chapter{总结与展望}

\section{全文工作总结}

本文围绕西南财经大学本科毕业论文排版规范要求，系统梳理了教务处对论文格式、字体字号、标题层级、页面边距、图表编排、公式编写、注释脚注与胶装顺序的全部技术指标，基于现代 \LaTeX{} 宏包技术研制了专属于西财本科生的毕业论文写作模板。

本文取得的主要成果包括：
\begin{enumerate}
    \item \textbf{版面规范 100\% 契合}：全文页边距严格设定为上下左右各 3 厘米，正文采用小四号宋体与 1.5 倍行距，一级标题采用二号华文中宋居中加粗，各级标题与正文间距均符合官方标准。
    \item \textbf{封面与声明全真复刻}：精确复现了西财本科专属单封面（含西南财经大学校名徽标、本科毕业论文（设计）题名与手写下划线信息栏），并提供符合规范的学术声明页与光华铁树封底。
    \item \textbf{现代化零依赖 JSON 配置}：全面支持通过 \texttt{config.json} 直观填写论文元数据，一键平滑切换终稿与匿名盲审模式，大幅降低了非计算机专业学生的使用门槛。
\end{enumerate}

\section{未来研究展望}

未来将进一步优化模板在 Overleaf 云端协同环境下的编译缓存性能，扩充面向经济学实证计量、金融工程数理推导以及量化回测图表的专用代码示例库，为更多西财学子提供坚实可靠的学术排版支持。
